\chapter{ISM} % ... + Lecture 09/03/2026

% missing previous lecture 

\section{ISM: Hot Ionized Medium}

The HIM is primarily created and energized by SN and the intense radiation from High-mass stars.

Hottest phase of the interstellar cycle, following the "Stardust injection” from stellar deaths.

HIM properties:

\begin{itemize}
    \item density: $\sim 10^{-3}$ cm$^{-3}$
    \item temperature: $\sim 10^6$ K — almost completely ionized gas
    \item thermal velocities: 100 km/s
    \item complex kinematics
    \item gas tracer: X-rays, high-ionization absorption lines: OVI, CIV, NV etc. (far UV)
    \item Occupies $\sim$50\% of the volume of the ISM
    \item The vertical scale height of this gas is $\sim$3 kpc, so it is sometimes referred to in the literature as the hot "corona" of the galaxy
    \item This hot gas is often buoyant and appears as bubbles, with characteristic dimensions of $\sim$20 pc, and fountains high above and below the disk.
    \item The gas is collisionally ionized and cools — primarily through adiabatic expansion and X-ray emission — on $\sim$Myr timescales
\end{itemize}

\section{ISM: Warm Ionized Medium}

Warm ionized medium (WIM) or Diffuse ionized gas (DIG) properties:

\begin{itemize}
    \item density: $\sim 0.2$ cm$^{-3}$
    \item temperature: $\sim 10^4$ K
    \item Ionized by photons escaping from HII regions
    \item Constitutes $\sim$90\% of all ionized hydrogen mass in the Galaxy ($\sim 10^9$ M$_\odot$)
    \item Fills $\sim$20\% of the ISM volume
    \item Transition Phase: It sits between the high-energy HIM and the cooler Diffuse Clouds
    \item Produces a faint, pervasive H$\alpha$ emission — red visible light that occurs when electrons recombine with hydrogen nuclei in "warm" gas — seen in every direction across the sky
    \item In other galaxies, observed as discrete loops filaments, and shell not directly associated with HII regions
    \item The DIG accounts for 30-50\% of the total H$\alpha$ emission of each galaxy 
    \item The H$\alpha$ intensity is measured across different frequency channels to understand the range of velocities caused by the Doppler effect.
\end{itemize}

\section{How do we observe the ISM}

The main idea is to decompose the observed emission into different wavelength bands.

\begin{itemize}
    \item \textbf{Radio continuum}: Synchrotron radiation from relativistic electrons (accelerated by SN shocks and pulsars) spiraling in the ~pG galactic magnetic field
    \item \textbf{Atomic hydrogen (21cm)}: HI neutral atomic hydrogen, hyperfine transition
    \item \textbf{Radio continuum (2.5 GHz)}: Synchrotron + bremsstrahlung (free-free) from supernova remanants and hot gas
    \item \textbf{Molecular hydrogen (CO)}: Traces molecular gas in the Galactic plane
    \item \textbf{Far Infrared}: Thermal emission from dust — heated by old stars (diffuse "cirrus") or by young
    massive stars in molecular clouds
    \item \textbf{Mid Infrared}: PAH emission lines from polycyclic aromatic hydrocarbons in the diffuse dust
    \item \textbf{Near Infrared}: Emission from old stars (main sequence + giants) with little dust absorption
    \item \textbf{Optical}: Emission from stars with dust absorption
    \item \textbf{X-ray}: X-ray binaries + diffuse hot gas; some absorption by neutral hydrogen
    \item \textbf{Gamma-ray}: Compact sources + gamma-ray photons from cosmic ray collisions with
    diffuse gas
\end{itemize}    

%todo: last slide